\begin{frame}[plain,c]\titlepage\end{frame}

\begin{frame}{Введение. Корреляционная функция.}
    \textbf{Корреляционная функция} (КФ) --- отношение измеренных двухчастичного инклюзивного и одночастичного инклюзивного спектров.
    \[C = \frac{A}{B} = \frac{A_{wei}}{A}\]
    \[A_{wei} = femtoweight \cdot A\]
    \[femtoweight = 1 + \cos\frac{(\Delta P\cdot\Delta r)}{hc}\]
    \centering, где $\Delta$ P --- разница 4-импульсов, а $\Delta$ r --- разница 4-векторов
\end{frame}

\begin{frame}{Введение. 1-мерная КФ.}
    В случае 1-мерной КФ используется инвариантный импульс пары:
    \[q_{inv} = \sqrt{{(P_{a} - P_{b})}^{2} - {(E_{a} - E_{b})}^{2}}\]
    Тогда КФ имеет следующий вид:
    \[C(q_{inv}) = \frac{A(q_{inv})}{B(q_{inv})} = \frac{A_{wei}(q_{inv})}{A(q_{inv})}\]
    А аппроксимирующая её кривая:
    \[C(q_{inv}) = 1 + \lambda\exp{\left(-\frac{R_{inv}^{2} q_{inv}^{2}}{(\hbar c)^2}\right)}\]
\end{frame}

\begin{frame}{Введение. 3-мерная КФ.}
    В случае 3-мерной КФ используется:
    \[q = p_{a} - p_{b}, \qquad P = p_{a} + p_{b}\]
    И система координат LCMS в которой любой 4-вектор $V$ из СЦМ представим в виде:

    \[V_{out } = (P_x V_x+P_y V_y) / P_T\]
    \[V_{side } = (P_x V_y-P_y V_x) / P_T\]
    \[V_{long} = (P_0 V_z-P_z V_0) / M_T\]
    \[\text{где } M_T^2=P_0^2-P_z^2 \text{ \;,a } P_T^2=P_x^2+P_y^2\]

\end{frame}

\begin{frame}{Введение. 3-мерная КФ.}
    Тогда 3-мерная КФ имеет вид:
    \[
        C =
        \frac{
            A(q_{out}, q_{side}, q_{long})
        }{
            B(q_{out}, q_{side}, q_{long})
        } =
        \frac{
            A_{wei}(q_{out}, q_{side}, q_{long})
        }{
            A(q_{out}, q_{side}, q_{long})
        }
    \]
    А аппроксимирующая её кривая:
    \[
        C(q) = 1 + \lambda \exp{(
            -\sum_{i,j\in\{o,s,l\}}
            \frac{
                R_{ij}q_{i}q_{j}
            }{(\hbar c)^2}
        )}
    \]
    В самом простом случае (без учета недиагональных элементов)
    \[
        C(q_{out}, q_{side}, q_{long}) =
        1 + \lambda \exp{(
            -\frac{
                R_{out}^{2} q_{out}^{2} +
                R_{side}^{2} q_{side}^{2} +
                R_{long}^{2} q_{long}^{2}
            }{(\hbar c)^2}
        )}
    \]
\end{frame}


\begin{frame}{Настройка модели и отбор частиц.}
	Для генерации использована транспортная модель SMASH:
    \begin{itemize}
        \item Столкновение:
            $^{197}_{79}\!Au + ^{197}_{79}\!Au$
        \item При энергии центра масс:
            $\sqrt{s_{NN}} = 3$ GeV
    \end{itemize}
    Для построения КФ смешивались частицы:
    \begin{itemize}
        \item Пионы (+) и (-).
        \item Произошло кинетическое вымораживание.
        \item $P_{total}$ $\in$ [0.15, 1.5] GeV/c
        \item $P_{T}$ $\in$ [0.15, 1.5] GeV/c
        \item $\eta$ $\in$ [-1.5, 1.5]
    \end{itemize}
    Также мы не учитываем события с упругим взаимодействием.
\end{frame}

\begin{frame}{Гистограммы событий.Прицельный параметр}
    \begin{figure}[H]
        \includegraphics[width=0.7\textwidth]{fig/impact_parameter.png}
        \caption{Распределене событий по прицельному параметру}\label{fig:sample}
    \end{figure}
\end{frame}

\setcounter{figure}{3}
\begin{frame}{Гистограммы частиц. 1.Координаты}
    \begin{figure}[H]
        \centering
        \begin{minipage}[b]{0.4\textwidth}
            \centering
            \includegraphics[width=\textwidth]{fig/x.png}
        \end{minipage}
        \begin{minipage}[b]{0.4\textwidth}
            \centering
            \includegraphics[width=\textwidth]{fig/y.png}
        \end{minipage}
        \begin{minipage}[b]{0.4\textwidth}
            \centering
            \includegraphics[width=\textwidth]{fig/z.png}
        \end{minipage}\label{fig:coordinates}
    \end{figure}
\end{frame}

\begin{frame}{Гистограммы частиц. 2.Время кинетического вымораживания}
    \begin{figure}[H]
        \includegraphics[width=0.7\textwidth]{fig/t.png}\label{fig:kineticfreezouttime}
        \caption{Распределение по времени наступления кинетического вымораживания}
    \end{figure}
\end{frame}

\setcounter{figure}{6}
\begin{frame}{Гистограммы частиц. 3.Импульсы}
    \begin{figure}[H]
        \centering
        \begin{minipage}[b]{0.4\textwidth}
            \centering
            \includegraphics[width=\textwidth]{fig/px.png}
        \end{minipage}
        \hfill
        \begin{minipage}[b]{0.4\textwidth}
            \centering
            \includegraphics[width=\textwidth]{fig/py.png}
        \end{minipage}
        \hfill
        \begin{minipage}[b]{0.4\textwidth}
            \centering
            \includegraphics[width=\textwidth]{fig/pz.png}
        \end{minipage}\label{fig:momentums}
    \end{figure}
\end{frame}

\begin{frame}{Гистограммы частиц. 4.Полный и поперечный импульсы}
    \begin{figure}[H]
        \centering
        \begin{minipage}[b]{0.47\textwidth}
            \centering
            \includegraphics[width=\textwidth]{fig/p_total.png}
        \end{minipage}
        \begin{minipage}[b]{0.47\textwidth}
            \centering
            \includegraphics[width=\textwidth]{fig/pt.png}
        \end{minipage}\label{fig:total_and_transvmomentums}
        \caption{Распределение по полному и поперечному импульсу для пионов}
    \end{figure}
\end{frame}

\begin{frame}{Гистограммы частиц. 5. Быстрота и псевдобыстрота}
    \begin{figure}[H]
        \centering
        \begin{minipage}[b]{0.47\textwidth}
            \centering
            \includegraphics[width=\textwidth]{fig/rapidity.png}
        \end{minipage}
        \begin{minipage}[b]{0.47\textwidth}
            \centering
            \includegraphics[width=\textwidth]{fig/pseudorapidity.png}
        \end{minipage}\label{fig:pseudo_and_rapidity}
        \caption{Распределение по быстроте и псевдобыстроте для пионов}
    \end{figure}
\end{frame}



\begin{frame}{Сравнение вычисленных параметров при разных центральностях}
    \begin{figure}[H]
        \centering
        \begin{minipage}[b]{0.47\textwidth}
            \centering
            \includegraphics[width=\textwidth]{fig/c_Rinv_vs_kt.png}
        \end{minipage}
        \begin{minipage}[b]{0.47\textwidth}
            \centering
            \includegraphics[width=\textwidth]{fig/c_lambda_vs_kt.png}
        \end{minipage}\label{fig:Rinv_lambda_4}
    \end{figure}

    $R_{inv}$ и $\lambda$ уменьшается с $k_{t}$ и от центральных к периферическим столкновениям
\end{frame}

\begin{frame}{1-мерная проекция 3-мерной КФ}
    \begin{figure}[H]
        \includegraphics[width=0.7\textwidth]{fig/c_1DCF_min.png}\label{fig:cf_3d_projection_1d}
    \end{figure}
\end{frame}

\begin{frame}{2-мерная проекция 3-мерной КФ}
    \begin{figure}[H]
        \includegraphics[width=0.7\textwidth]{fig/c_2DCF_min.png}\label{fig:cf_3d_projection_2d}
    \end{figure}
\end{frame}

\begin{frame}{Сравнение вычисленных параметров из 3-мерной КФ при разны центральностях}
    \begin{figure}[H]
        \includegraphics[width=0.65\textwidth]{fig/c_out_side_long_lambda.png}\label{fig:radii_lambda_centrality_3d}
    \end{figure}
\end{frame}

\section{Статистика корреляционной функции}
\begin{frame}{КФ как средний вес пары}
    \small
    $A(\mathbf q)$ --- распределение пар в числителе,
    $B(\mathbf q)$ --- опорное распределение; $C=\mathcal{N}A/B$.
    $\mathcal{N}$ --- нормировка, $n$ --- число пар в ячейке.
    Здесь $\mathcal{N}=1$, обе гистограммы заполняются по одним $n$ парам:
    \[
        A = \sum_{i=1}^{n} w_i, \qquad
        B = \sum_{i=1}^{n} 1 = n, \qquad Q_A = \sum_{i=1}^{n} w_i^2.
    \]
    Поэтому КФ в одной ячейке представляет собой средний вес:
    \[
        C = \frac{A}{B} = \frac{A}{n} = \overline{w}, \qquad n>0.
    \]
    \begin{itemize}
        \item $Q_A$ хранится в \texttt{Sumw2} взвешенной гистограммы.
        \item При случайном $n$ общие пары могут связывать суммы $A$ и $B$.
    \end{itemize}
\end{frame}

\begin{frame}{Ошибка среднего при фиксированном числе пар}
    \small
    Пусть $n$ фиксировано, а веса независимы и одинаково распределены
    с дисперсией $\sigma_w^2$. Тогда
    \[
        \operatorname{Var}(C\mid n) = \frac{\sigma_w^2}{n}.
    \]
    Несмещённая выборочная оценка при $n>1$:
    \[
        s_w^2 = \frac{Q_A-A^2/n}{n-1}, \qquad
        \widehat{\operatorname{Var}}(C\mid n)
        = \frac{s_w^2}{n}
        = \frac{Q_A-A^2/n}{n(n-1)}.
    \]
    \begin{itemize}
        \item $\sigma_w^2$ --- дисперсия распределения весов;
              $s_w^2$ --- её оценка по наблюдаемым парам.
        \item Поправка $n-1$ учитывает оценивание среднего по той же выборке.
    \end{itemize}
\end{frame}

\begin{frame}{Частный случай: бинарный отбор}
    \small
    Для индикатора отбора $I_i\in\{0,1\}$ выполнено $I_i^2=I_i$.
    Поэтому $Q_A=A$, а $\widehat p=A/B=A/n$ оценивает вероятность отбора $p$.
    \[
        \operatorname{Var}(\widehat p\mid n) = \frac{p(1-p)}{n}.
    \]
    Две оценки этой дисперсии:
    \[
        \begin{aligned}
            \widehat V_{\mathrm{ROOT\ B}}
                &= \frac{\widehat p(1-\widehat p)}{n},
                &&\text{подстановка }p\to\widehat p;\\
            \widehat V_{\texttt{cf\_maker}}
                &= \frac{\widehat p(1-\widehat p)}{n-1},
                &&\text{несмещённая оценка},\quad n>1.
        \end{aligned}
    \]
    Одного ограничения $0\leq w_i\leq 1$ недостаточно:
    для небинарных весов равенство $Q_A=A$ обычно не выполняется.
\end{frame}

\begin{frame}{Общие суммы: пуассоновская модель}
    \small
    Здесь число вкладов случайно: $n\sim\operatorname{Pois}(\mu)$.
    Вклады $(a_i,b_i)$ независимы, одинаково распределены и не зависят от $n$.
    \[
        A=\sum_i a_i, \qquad B=\sum_i b_i, \qquad C=\frac{A}{B}.
    \]
    Суммы для оценки дисперсий и ковариации:
    \[
        Q_A=\sum_i a_i^2, \qquad
        Q_B=\sum_i b_i^2, \qquad Q_{AB}=\sum_i a_i b_i.
    \]
    $\widehat V_A=Q_A$, $\widehat V_B=Q_B$,
    $\widehat{\operatorname{Cov}}(A,B)=Q_{AB}$.
    Разложение отношения до первого порядка даёт
    \[
        \widehat{\operatorname{Var}}(C)\simeq
        \frac{Q_A-2C Q_{AB}+C^2 Q_B}{B^2},
        \qquad B\ne 0.
    \]
    Это дельта-оценка, а не точная дисперсия отношения для малой выборки.
\end{frame}

\begin{frame}{Ковариация взвешенной суммы и числа пар}
    \small
    Для КФ в пуассоновской модели $a_i=w_i$, $b_i=1$. Поэтому
    \[
        \widehat V_A=Q_A, \qquad \widehat V_B=n, \qquad
        \widehat{\operatorname{Cov}}(A,B)=\sum_i w_i\cdot 1=A.
    \]
    Подставляя $C=A/B$ и $B=n$ в дельта-оценку, получаем
    \[
        \widehat V_{\Delta}=
        \frac{Q_A-2CA+C^2B}{B^2}
        =\frac{Q_A-A^2/n}{n^2}.
    \]
    Оценка \texttt{cf\_maker} для фиксированного $n$ отличается поправкой:
    \[
        \widehat{\operatorname{Var}}(C\mid n)
        =\frac{n}{n-1}\,\widehat V_{\Delta}, \qquad n>1.
    \]
    При большом $n$ оценки асимптотически совпадают.
    Их статистические допущения различаются.
\end{frame}

\begin{frame}{КФ и взвешенная эффективность: разные вклады}
    \small
    Для общих пар различается связь числителя и знаменателя:
    \[
        \begin{array}{l|c|c|c}
            & a_i & b_i & \widehat{\operatorname{Cov}}(A,B)\\ \hline
            \text{КФ} & w_i & 1 & A\\
            \text{эффективность} & u_i I_i & u_i & Q_{\mathrm{pass}}
        \end{array}
    \]
    Для эффективности $\varepsilon=A/B$, $I_i\in\{0,1\}$, $u_i\geq 0$ и
    \[
        Q_{\mathrm{pass}}=\sum_i u_i^2 I_i=\widehat V_A,
        \qquad Q_{\mathrm{all}}=\sum_i u_i^2.
    \]
    Дельта-оценка эффективности имеет вид формулы ROOT \texttt{B}:
    \[
        \widehat{\operatorname{Var}}(\varepsilon)\simeq
        \frac{(1-2\varepsilon)Q_{\mathrm{pass}}+\varepsilon^2 Q_{\mathrm{all}}}{B^2}.
    \]
    Произвольные веса $w_i$ в числителе не обеспечивают эту ковариацию.
    \par\medskip
    {\footnotesize Источник:
        \href{https://root.cern.ch/doc/master/TH1_8cxx_source.html}{ROOT: TH1::Divide, опция B}.}
\end{frame}

\begin{frame}{Контрольный пример: произвольные веса}
    \small
    По $n=100$ пар, $B=n$. Проверка: ROOT 6.40.02.
    \begin{center}
        \renewcommand{\arraystretch}{1.35}
        \begin{tabular}{lccc}
            \hline
            Веса & $C$ & $\sigma_C$, \texttt{cf\_maker} & $\sigma_C$, ROOT \texttt{B}\\ \hline
            Все $w_i=1$ & 1 & 0 & 0\\
            50 весов 0, 50 весов 2 & 1 & 0.1005 & 0\\
            Все $w_i=2$ & 2 & 0 & 0.2828\\ \hline
        \end{tabular}
        \par\smallskip
        Контрольный пример, не данные столкновений.
    \end{center}
    \begin{itemize}
        \item Во второй строке $A=B$, но веса изменяются, и ошибка среднего ненулевая.
        \item ROOT \texttt{B} задаёт нулевую ошибку при $A=B$.
              Это не проверка применимости модели эффективности.
    \end{itemize}
\end{frame}

\begin{frame}{Корреляции пар и применимость ошибок}
    \small
    Число пар $n$ фиксировано, но веса могут быть зависимы:
    \[
        \operatorname{Var}(C\mid n)=\frac{1}{n^2}
        \left[\sum_i\operatorname{Var}(w_i\mid n)
        +2\sum_{i<j}\operatorname{Cov}(w_i,w_j\mid n)\right].
    \]
    \begin{itemize}
        \item Общие события или треки могут связывать веса разных пар.
        \item Bootstrap или jackknife по событиям сохраняют
              совместные вклады пар внутри события.
        \item При $n\leq 1$ выборочная дисперсия недоступна;
              нулевая записанная ошибка не означает точное измерение.
    \end{itemize}
    Опция ROOT \texttt{B} не устраняет событийные корреляции.
    Для интерпретации $\chi^2$ нужно проверить допущения об ошибках.
\end{frame}


\section{Обновление программного анализа}

\begin{frame}{Обновление анализа. Точность исходных гистограмм}
    \small
    \begin{itemize}
        \item Добавлено чтение исходных \texttt{TH3F} и \texttt{TH3D} с сохранением
              моментов, осей и числа заполнений.
        \item Преобразование \texttt{TH3F} в \texttt{TH3D} сохраняет данные,
              но не восстанавливает потерянную при накоплении точность.
        \item Если $Q-S^2/N$ неразличимо на фоне округления, значение $S/N$
              сохраняется. Ошибка такого бина недоступна, и он исключается из
              $\chi^2$-аппроксимации. При $N=1$ ошибка также недоступна;
              пустые бины не оцениваются.
        \item Нарушения моментов за пределами оценки округления приводят к
              явной ошибке с диагностикой $N$, $S$ и $Q$.
    \end{itemize}
    Для следующих генераций требуется накопление исходных сумм в двойной точности.
\end{frame}

\begin{frame}{Обновление анализа. Проекции и сравнение зарядов}
    \small
    Проекция КФ учитывает число пар в ячейках среза $\mathcal{W}$:
    \[
        C_{\mathrm{proj}}(q_i) =
        \frac{\sum_{j\in\mathcal{W}} S_j}{\sum_{j\in\mathcal{W}} N_j}.
    \]
    Отношение проекций и среднее отношения по $K$ допустимым ячейкам:
    \[
        \frac{C^{++}_{\mathrm{proj}}}{C^{--}_{\mathrm{proj}}},
        \qquad
        \left\langle\frac{C^{++}}{C^{--}}\right\rangle_{\mathcal{W}}
        = \frac{1}{K}\sum_{j\in\mathcal{W}}\frac{C^{++}_j}{C^{--}_j}.
    \]
    \begin{itemize}
        \item Среднее отношения одинаковых ненулевых КФ равно 1 при $K>0$.
        \item В 2-мерных проекциях первая выбранная LCMS-ось соответствует $X$,
              вторая --- $Y$.
        \item Бины среза выбираются по фактическим границам с учётом округления.
    \end{itemize}
\end{frame}

\begin{frame}{Обновление анализа. Контроль аппроксимации}
    \small
    \begin{itemize}
        \item В графики параметров и модельные кривые включаются успешные
              конечные результаты с $\mathrm{ndf}>0$ и без попадания
              свободных параметров на границы.
        \item При повторной аппроксимации пригодный результат имеет приоритет
              перед граничным. $\chi^2$ сравнивается внутри одного класса
              успешных результатов.
        \item Сохраняются параметры, ошибки, корреляции, статусы,
              $\chi^2/\mathrm{ndf}$ и число фактических попыток.
        \item Повторный запуск этапов использует сохранённые результаты только
              при совпадении исходного файла и настроек анализа.
        \item Непригодные фиты сохраняются для диагностики и дают ненулевой
              код завершения с указанием причин.
    \end{itemize}
\end{frame}

\begin{frame}{Обновление анализа. Проверка результатов}
    \small
    \begin{itemize}
        \item Локально пройдены все 14 проверок CTest в обычной сборке
              и со средствами ASan/UBSan (macOS, ROOT 6.40).
        \item На контрольном \texttt{TH3F}-входе проверены все этапы анализа,
              восстановление заданных параметров, ROOT-выходы и изображения.
        \item По логам от 4 октября 2026 г. запуск по быстроте
              \texttt{y\_run\_3609\_fpc} завершил все этапы:
              80 из 80 фитов прошли критерии пригодности.
        \item В запуске по $k_t$ \texttt{kt\_run\_2707\_fpc} сохранена диагностика
              6 непригодных фитов из 32. После исправления выбора повторного
              фита этот набор ещё требует проверки.
    \end{itemize}
    Прохождение критериев пригодности не заменяет проверку физических
    допущений и статистической модели. Прежние рисунки выше не пересчитаны
    в рамках этого обновления.
\end{frame}

% XeLaTeX searches the output directory first, so each build reads its own chapter.
\IfFileExists{\jobname-results.tex}{\input{\jobname-results.tex}}{}

\begin{frame}{Заключение}
    \begin{itemize}
	\item Был сгенерирован набор частиц в транспортной модели “SMASH”.
        \item Были построены и изучены гистограммы для базового анализа столкновения.
        \item На базе сгенерированных модельных данных был освоен метод фемтоскопических корреляций.
        \item Были построены корреляционные функции для одномерного и трехмерного анализа.
        \item Были получены зависимости радиусов и силы корреляции от поперечного импульса и центральности.
        \item Был обновлён и проверен программный анализ: статистические ошибки,
              проекции, сравнение зарядов и контроль качества аппроксимации.
    \end{itemize}
\end{frame}

\begin{frame}{Цели}
    \begin{itemize}
	\item  Набрать более полную статистику для высокой центральности и для разных значений $\sqrt{S_{NN}}$.
        \item Изучить асимметрию. В частности для направления out.
        \item Проверить соответствие модельных и экспериментальных данных.
    \end{itemize}
\end{frame}

\begin{frame}
	\Large \textbf{Спасибо за внимание!}
\end{frame}

\begin{frame}{Дополнительные слайды. КФ для всех центральностей и $k_{t}$ классов}
    \begin{figure}[H]
        \includegraphics[width=0.65\textwidth]{fig/c_all_1d_CF.png}\label{fig:all_cf_1d}
    \end{figure}
\end{frame}

\begin{frame}{Дополнительные слайды. 3-мерные КФ для всех центральностей и $k_{t}$ классов}
    \begin{figure}[H]
        \includegraphics[width=0.65\textwidth]{fig/c_all_3d_CF.png}\label{fig:all_cf_3d}
    \end{figure}
\end{frame}

\begin{frame}{Дополнительные слайды. 1-мерные проекции 3-мерных КФ для всех центральностей и $k_{t}$ классов, на ось out}
    \begin{figure}[H]
        \includegraphics[width=0.65\textwidth]{fig/c_all_3d_CF_proj_out.png}\label{fig:all_cf_3d_projection_out}
    \end{figure}
\end{frame}

\begin{frame}{Дополнительные слайды. 1-мерные проекции 3-мерных КФ для всех центральностей и $k_{t}$ классов, на ось side}
    \begin{figure}[H]
        \includegraphics[width=0.65\textwidth]{fig/c_all_3d_CF_proj_side.png}\label{fig:all_cf_3d_projection_side}
    \end{figure}
\end{frame}

\begin{frame}{Дополнительные слайды. 1-мерные проекции 3-мерных КФ для всех центральностей и $k_{t}$ классов, на ось long}
    \begin{figure}[H]
        \includegraphics[width=0.65\textwidth]{fig/c_all_3d_CF_proj_long.png}\label{fig:all_cf_3d_projection_long}
    \end{figure}
\end{frame}

\begin{frame}{Дополнительные слайды. 2-мерные проекции 3-мерных КФ для всех центральностей и $k_{t}$ классов, в осях out-side}
    \begin{figure}[H]
        \includegraphics[width=0.65\textwidth]{fig/c_2DCF_out_vs_side.png}\label{fig:all_cf_3d_projection_out_side}
    \end{figure}
\end{frame}

\begin{frame}{Дополнительные слайды. 2-мерные проекции 3-мерных КФ для всех центральностей и $k_{t}$ классов, в осях out-long}
    \begin{figure}[H]
        \includegraphics[width=0.65\textwidth]{fig/c_2DCF_out_vs_long.png}\label{fig:all_cf_3d_projection_out_long}
    \end{figure}
\end{frame}

\begin{frame}{Дополнительные слайды. 2-мерные проекции  3-мерных КФ для всех центральностей и $k_{t}$ классов, в осях side-long}
    \begin{figure}[H]
        \includegraphics[width=0.65\textwidth]{fig/c_2DCF_side_vs_long.png}\label{fig:all_cf_3d_projection_side_long}
    \end{figure}
\end{frame}
